% Figure: a square wave reconstructed as a sum of harmonic phasors. The
% left panel shows the partial sums N = 1, 3, 5 of the Fourier series for
% a unit square wave; the right panel shows the harmonic amplitudes (the
% magnitude spectrum) as a stem plot, each stem the magnitude of one
% harmonic phasor.
\begin{figure}[htbp]
\centering
\begin{tikzpicture}[scale=1.0, line cap=round, >=latex]
  % Left panel: partial sums.
  \begin{scope}
    \draw[->, engblack, thin] (-0.2, 0) -- (5.0, 0) node[right, font=\small] {$\omega t$};
    \draw[->, engblack, thin] (0, -1.7) -- (0, 1.7) node[above, font=\small] {$x(t)$};
    % Ideal square wave (light), period 2pi mapped to 0..4.8 (i.e. ~1.5 periods).
    \draw[engblack!35, thick]
      (0,1) -- (2.4,1) -- (2.4,-1) -- (4.8,-1);
    % N = 1 (pure fundamental): (4/pi) sin. Amplitude ~1.27.
    \draw[engblue, thin, domain=0:4.8, samples=80, smooth]
      plot (\x, {1.273*sin(deg(\x*pi/2.4))});
    % N = 3 (fundamental + third).
    \draw[engred, thick, domain=0:4.8, samples=160, smooth]
      plot (\x, {1.273*sin(deg(\x*pi/2.4)) + 0.424*sin(deg(3*\x*pi/2.4))});
    % N = 5.
    \draw[engamber, very thick, domain=0:4.8, samples=240, smooth]
      plot (\x, {1.273*sin(deg(\x*pi/2.4)) + 0.424*sin(deg(3*\x*pi/2.4)) + 0.255*sin(deg(5*\x*pi/2.4))});
    \node[engblue, font=\footnotesize, anchor=south] at (1.2, 1.05) {$N{=}1$};
    \node[engamber, font=\footnotesize, anchor=south west] at (3.1, 0.55) {$N{=}5$};
    \node[font=\small, anchor=north] at (2.4, -1.85) {(a) partial sums};
  \end{scope}
  % Right panel: stem plot of harmonic magnitudes.
  \begin{scope}[xshift=6.6cm]
    \draw[->, engblack, thin] (-0.2, 0) -- (4.4, 0) node[right, font=\small] {harmonic $n$};
    \draw[->, engblack, thin] (0, -0.2) -- (0, 1.9) node[above, font=\small] {$|\tilde X_n|$};
    % Stems at n = 1,3,5,7,9 with amplitudes 4/(n pi).
    \foreach \n/\h in {1/1.6, 3/0.533, 5/0.32, 7/0.229, 9/0.178} {
      \draw[engblue, very thick] ({\n*0.42}, 0) -- ({\n*0.42}, \h);
      \filldraw[engblue] ({\n*0.42}, \h) circle (0.045);
      \node[engblack, font=\scriptsize, anchor=north] at ({\n*0.42}, 0) {$\n$};
    }
    \node[font=\small, anchor=north] at (2.1, -0.55) {(b) magnitude spectrum};
  \end{scope}
\end{tikzpicture}
\caption[A square wave as a sum of harmonic phasors]{A unit square wave
reconstructed from its Fourier series. Panel (a): the partial sums
$N = 1$ (fundamental alone), $N = 3$, and $N = 5$ converge toward the
ideal square wave (grey). Each harmonic is a phasor at frequency
$n\omega$; the time-domain wave is the real part of the sum
$\sum_n \tilde X_n e^{in\omega t}$. Panel (b): the magnitude spectrum,
the magnitudes $|\tilde X_n| = 4/(n\pi)$ of the odd-harmonic phasors,
drawn as a stem plot. The even harmonics are absent by the wave's
half-wave symmetry. The overshoot near each discontinuity in panel (a)
is the Gibbs phenomenon, which does not vanish as $N$ grows.}
\label{fig:vol02:ch03:fourier-square}
\end{figure}
